\section{Introduction}
\label{sec:introduction}

VoltShare coordinates a network of electric-vehicle charging stations. Drivers find a
station, reserve a connector, plug in and are billed; the operator sees the network from a
back office. Underneath, the system decides three things: which vehicle gets which connector,
how the limited electrical capacity of a site is shared among the cars charging at it, and
what happens to both when a node fails.

This domain represents a real contention scenario: different drivers could request the same connector at the same time. This problem cannot be resolved in an easy way by a central authority that assigns arbitrarly the driver
to a resource, because they have their own destination and can ignore any suggestion.
The system can only \emph{arbitrate the requests that drivers make}.


\subsection{The problems this project solves}

Four of them:

\begin{itemize}
  \item \textbf{Exclusive access to a physical resource.} Several drivers ask for the same
        connector at the same instant, and the outlet must end up promised to one of them.
  \item \textbf{One reservation per vehicle across the whole network.} This is the hard one.
        Two stations can each be locally correct and still break the guarantee together, so
        it needs an authority, which then needs replication, election and a quorum to
        survive its own failure.
  \item \textbf{Sharing a divisible resource.} Tipically, in a real word scenario, the total site capacity is lower than the sum of what its
        connectors are rated for, so the available power is apportioned among active
        sessions and renegotiated whenever a car arrives, leaves or fills up.
  \item \textbf{Failures that must not take resources with them.} Nodes crash and links
        break, and from the outside the two are indistinguishable. A reservation held by a
        node that has stopped answering must be released without waiting for it to return.

\end{itemize}

\Cref{sec:problems} states all seven problems precisely, \cref{sec:addressing} gives the
mechanism chosen for each, and \cref{sec:fault-tolerance} covers the failures.

The outlet, its contactor, its meter
and the car are an emulator speaking a reduced set of OCPP 1.6-J messages (see \cref{sec:api-ws-cp}), so the boundary
falls where a real interface already exists (\cref{sec:boundary}). Seven containers on one
host run everything else (\cref{sec:deployment}).

Measurements come from running deployments and where something was not measured, the text says so,
and the limits are declared in \cref{sec:no-tolerance}.
